\chapter{Розпізнавання зображень і відео}
\chaptermark{Розпізнавання}
\label{sec:recognition}

\section{Задача: те саме за різних пікселів}

Розділ~\ref{sec:sensors} довів зображення до входу машини: близько мільйона вихідних нейронів ока, стиснений потік, у якому передаються здебільшого краї й зміни. Але між датчиком і дією є крок, про який ми досі мовчали. Кішка на картинці --- це не набір пікселів. Зсуньте її на пів кадру, поверніть, віддаліть, вимкніть половину світла --- майже всі пікселі стануть іншими, а відповідь «кішка» має лишитися тією самою. Інженер називає це \emph{інваріантністю}: відповідь класифікатора не повинна змінюватися від зсуву, масштабу, повороту й освітлення.

Труднощі добре видно на простому досліді. Візьмімо одновимірну «сітківку» з $24$ точок і дві фігури по п'ять точок, які можуть стояти будь-де. Класифікатор, що порівнює вхід зі зразком, запам'ятаним в одному місці, вгадує в $49$--$52\%$ випадків --- стільки ж, скільки монетка. Зсув цілком ламає порівняння за пікселями. Потрібна інша схема.

\section{Конвеєр зі ступенів}

Як швидко мозок це робить, ми вже знаємо з розділу~\ref{sec:code}: тварину на миттєво показаній картинці впізнають приблизно за $150\ms$, і сигнал проходить близько десятка ступенів, по $10$--$20\ms$ на кожен~\cite{Thorpe1996}. На кожному ступені нейрон встигає видати нуль або один імпульс. Отже, швидке розпізнавання --- це один прохід конвеєром, без довгих роздумів і повторів. Питання в тому, що робить кожен ступінь.

Відповідь, яку підказали вимірювання на перших ступенях зору~\cite{HubelWiesel1962}, складається з двох почергових операцій (рис.~\ref{fig:conveyor}).

\emph{Вибірковість.} Нейрон ступеня $S$ дивиться на невелику ділянку входу й спрацьовує, лише якщо там є певне поєднання частин: ось цей край, а поруч той. Це І за частинами --- пороговий нейрон розділу~\ref{sec:glutcomputer}, поріг якого вищий, ніж дає будь-яка одна частина.

\emph{Інваріантність.} Нейрон ступеня $C$ збирає виходи багатьох $S$-нейронів, що шукають те саме поєднання в різних місцях, і спрацьовує, якщо воно знайшлося хоч десь. Це АБО за положеннями --- а точніше, вибір найбільшого.

Для одновимірної моделі обидві операції записуються так:
\begin{empheq}[box=\keybox]{equation}
s_{f,p} \;=\; \Bigl[\sum_i t_{f,i}\,x_{p+i} - \theta\Bigr]_+ ,
\qquad
c_f \;=\; \max_p\, s_{f,p} ,
\label{eq:scstage}
\end{empheq}
де $x$ --- вхід, $t_{f,i}$ --- зразок ознаки $f$, $p$ --- положення, $\theta$ --- поріг. Перший вираз робить відповідь вибірковою, другий --- незалежною від положення. Найбільше з багатьох входів мережа отримує засобами розділу~\ref{sec:gaba}: взаємне гальмування залишає найсильніший вхід (твердження~\ref{prop:wta}), а ділення на загальну активність вирівнює відповіді за різної яскравості~\cite{Carandini2012}.

\begin{figure}[t]
\centering
\begin{tikzpicture}[>=Stealth,
  px/.style={draw,minimum size=0.36cm,inner sep=0pt},
  on/.style={px,fill=blue!55!black},
  snode/.style={circle,draw,very thick,minimum size=0.62cm,inner sep=0pt,font=\scriptsize,fill=blue!6},
  cnode/.style={circle,draw,very thick,minimum size=0.8cm,inner sep=0pt,font=\scriptsize,fill=red!10}]
% two inputs: the same shape at two positions
\foreach \row/\sh/\lab in {0/1/{кадр 1},1/7/{кадр 2}}{
  \begin{scope}[yshift=-\row*1.0cm]
  \node[left,font=\scriptsize] at (-0.25,0) {\lab};
  \foreach \i in {0,...,13}{\node[px] at (\i*0.4,0) {};}
  \foreach \d in {0,1,3,4}{\node[on] at ({(\sh+\d)*0.4},0) {};}
  \end{scope}}
% S stage: detectors of the shape at several positions
\foreach \k in {0,...,5}{
  \node[snode] (S\k) at ({0.8+0.8*\k},1.7) {$S$};}
\node[font=\scriptsize,left] at (-0.25,1.7) {ступінь $S$: І за частинами};
\draw[->,thick,blue!60!black] (1.2,0.25) -- (S1.south);
\draw[->,thick,orange!85!black] (3.6,0.25) -- (S4.south);
\node[font=\scriptsize,blue!60!black,anchor=east] at (1.25,0.85) {кадр 1};
\node[font=\scriptsize,orange!85!black,anchor=west] at (3.9,0.85) {кадр 2};
% C stage
\node[cnode] (C) at (2.8,3.25) {$C$};
\node[font=\scriptsize,left] at (-0.25,3.25) {ступінь $C$: АБО за положеннями};
\foreach \k in {0,...,5}{\draw[->,gray!60!black] (S\k.north) -- (C);}
\draw[->,very thick,red!70!black] (C.north) -- ++(0,0.6) node[above,font=\scriptsize,black] {«фігура є» --- в обох кадрах};
% next stages
\draw[decorate,decoration={brace,amplitude=5pt},gray!60!black] (6.3,1.4) -- (6.3,-1.3);
\node[right,font=\scriptsize,align=left] at (6.5,0.05) {той самий\\об'єкт, інші\\пікселі};
\draw[->,thick,densely dashed,gray!60!black] (3.6,3.25) -- (5.9,3.25) node[right,font=\scriptsize,black,align=left] {наступні пари\\$S$, $C$ --- більші,\\складніші, інваріантніші};
\end{tikzpicture}
\caption{Одна пара ступенів конвеєра. Нейрони $S$ шукають те саме поєднання частин у різних місцях; нейрон $C$ відповідає, якщо воно знайшлося хоч десь~\eqref{eq:scstage}. Фігура в кадрі 1 і в кадрі 2 збуджує різні $S$-нейрони, але той самий $C$. Наступні пари ступенів повторюють те саме на дедалі більших і складніших поєднаннях.}
\label{fig:conveyor}
\end{figure}

У тій самій одновимірній моделі пара ступенів~\eqref{eq:scstage} упізнає фігуру в будь-якому місці без помилок, а якщо кожну точку входу випадково перевернути з імовірністю $0.1$ --- у $86\%$ випадків, з імовірністю $0.2$ --- у $70\%$. Монетка, як і раніше, дає половину. Покладіть кілька таких пар одну над одною: перша шукає краї, наступна --- поєднання країв, ще вище --- частини предметів, нагорі --- предмети цілком. З кожною парою поєднання більші, а відповідь дедалі менше залежить від того, де і як лежить предмет~\cite{Riesenhuber1999}.

\begin{keyblock}
\begin{proposition}[Конвеєр вибірковості й інваріантності]
\label{prop:conveyor}
Швидке розпізнавання --- один прохід десятком ступенів. На кожній парі ступенів І за частинами~\eqref{eq:scstage} робить відповідь вибірковою, а АБО за положеннями --- незалежною від зсуву. Чергування цих двох операцій дає відповідь, яка розрізняє предмети й майже не залежить від того, де і як їх видно. Будову перших ступенів і швидкість проходу виміряно; те, що весь ланцюжок зводиться до цих двох операцій, --- спрощення.
\end{proposition}
\end{keyblock}

Якщо ця будова здається знайомою, так і має бути. Саме її чергування --- пошук ознаки в усіх місцях, потім об'єднання за сусідніми місцями --- інженери перенесли в штучні \emph{згорткові} мережі~\cite{Fukushima1980}. А нещодавно пішли у зворотному напрямку: згорткові мережі, навчені розпізнавати предмети, виявилися найкращою з відомих моделей відповідей нейронів на верхніх ступенях зору~\cite{Yamins2014}. Результат нагорі читається просто: за сумарними частотами близько сотні нейронів останнього ступеня лінійний вирішальний блок упізнає предмет через десяту частку секунди після показу~\cite{Hung2005}. Це частотний код групи з розділу~\ref{sec:code}.

\section{Учитель, якого немає: вчитися на відео}

Згорткову мережу навчають на мільйонах картинок із підписами. Мозку підписів майже ніхто не дає: дитина бачить предмети годинами, а слово «чашка» чує зрідка. Звідки тоді ступені $C$ знають, які $S$-нейрони об'єднувати? Адже щоб об'єднати «цю фігуру тут» із «цією фігурою там», треба вже знати, що це та сама фігура.

Підказку дає саме відео. Світ змінюється неперервно: предмет у полі зору лишається тим самим предметом, доки він зсувається, повертається й змінює освітлення, і лише зрідка погляд переходить на інший. Усе, що змінюється \emph{повільно}, найімовірніше, належить до предмета, а все, що змінюється швидко, --- до його положення й вигляду~\cite{WiskottSejnowski2002}. Отже, $C$-нейронові досить правила Гебба з розділу~\ref{sec:dofa} зі слідом: пов'язувати те, що йшло поспіль, а не лише те, що йшло одночасно~\cite{Foldiak1991}:
\begin{empheq}[box=\keybox]{equation}
\tau_{\text{сл}}\,\frac{\dd\bar y_k}{\dd t} \;=\; -\bar y_k + y_k ,
\qquad
\frac{\dd\mathbf w_k}{\dd t} \;=\; \eta\,\bar y_k\,\bigl(\mathbf x - \mathbf w_k\bigr) .
\label{eq:traceinvariance}
\end{empheq}
Тут $y_k$ --- активність $C$-нейрона $k$, який виграв змагання через гальмування, $\bar y_k$ --- її слід, те саме RC-коло, що й у правилі~\eqref{eq:threefactor}, а $\mathbf x$ --- виходи $S$-нейронів. Нейрон, що виграв на першому вигляді предмета, ще пам'ятає це, коли надходить другий вигляд, і підтягує до себе й його.

\begin{figure}[t]
\centering
\begin{tikzpicture}[>=Stealth,x=1.0cm,y=1.0cm]
\foreach \i/\lab in {0/1,1/2,2/3,3/4}{
  \node[draw,thick,fill=blue!8,minimum width=0.9cm,minimum height=0.55cm,font=\scriptsize] at (\i+0.5,2.2) {$A$ у \lab};}
\foreach \i/\lab in {0/4,1/3,2/2,3/1}{
  \node[draw,thick,fill=orange!12,minimum width=0.9cm,minimum height=0.55cm,font=\scriptsize] at (\i+6.5,2.2) {$B$ у \lab};}
\node[font=\scriptsize,gray!40!black] at (5.0,2.2) {пауза};
\draw[->,gray!70!black] (0,0) -- (10.4,0) node[below left,font=\scriptsize,black] {час};
% traces
\draw[thick,blue!60!black] (0,0.05) .. controls (0.4,1.2) and (0.8,1.3) .. (1.2,1.35) -- (4,1.4) .. controls (4.6,0.5) and (5.2,0.15) .. (6,0.08) -- (10,0.05);
\draw[thick,orange!85!black] (0,0.05) -- (6,0.05) .. controls (6.4,1.2) and (6.8,1.3) .. (7.2,1.35) -- (10,1.4);
\node[font=\scriptsize,blue!60!black,anchor=west] at (1.4,1.65) {слід $\bar y_1$ --- тримається весь прохід $A$};
\node[font=\scriptsize,orange!85!black,anchor=west] at (7.2,1.65) {слід $\bar y_2$};
\draw[decorate,decoration={brace,amplitude=4pt},gray!60!black] (4.0,2.6) -- (0,2.6);
\draw[decorate,decoration={brace,amplitude=4pt,mirror},gray!60!black] (0,-0.35) -- (4,-0.35) node[midway,below=4pt,font=\scriptsize,black] {усі вигляди $A$ пов'язуються з нейроном 1};
\end{tikzpicture}
\caption{Як відео навчає інваріантності, схема. Предмет $A$ проходить через чотири положення, потім пауза, потім предмет $B$. Слід~\eqref{eq:traceinvariance} нейрона, що виграв на першому вигляді $A$, тримається весь прохід, і всі вигляди $A$ виявляються пов'язаними з одним нейроном. Пауза стирає слід, і $B$ дістається іншому нейронові.}
\label{fig:tracevideo}
\end{figure}

Ми перевірили це на моделі (рис.~\ref{fig:tracevideo}). Два $C$-нейрони змагаються за входи від $16$ $S$-нейронів: дві фігури у восьми положеннях. Фігура проїжджає через усі положення, по $50\ms$ на кожне, потім пауза $0.3\,\text{с}$ і наступна випадкова фігура. Жодних підписів. Якщо слід короткий, $\tau_{\text{сл}}=5\ms$, нейрони ділять вхід за \emph{положенням}, і відповідь збігається з фігурою не краще, ніж у монетки: $52\%$ у середньому за десять прогонів. Якщо слід порівнянний із часом показу одного положення, від $\tau_{\text{сл}}\approx20\ms$, кожен нейрон збирає \emph{всі} положення своєї фігури (за $20$, $50$ і $200\ms$ --- в усіх десяти прогонах): інваріантність вивчено на самому лише відео. Пауза між предметами важлива: без неї слід перетікає на наступний предмет і плутає їх.

Те саме видно й у досліді. Якщо в експерименті непомітно підміняти предмет, доки око перескакує з одного місця на інше, то нейрони верхніх ступенів за годину-дві починають відповідати на підмінений предмет як на той самий: вони пов'язують те, що йшло поспіль~\cite{LiDiCarlo2008}. Інваріантність у мозку справді вчиться на неперервності часу. Наскільки цим правилом пояснюється все навчання зору, --- гіпотеза.

\section{Відео --- це рух}

Відео --- не лише потік кадрів для навчання, а й сама задача: що рухається, куди і як швидко. Перший крок нам знайомий --- детектор напрямку розділу~\ref{sec:gaba} (рис.~\ref{fig:direction}), у якому запізніле гальмування гасить рух в один бік. Такі детектори стоять по всьому полю зору, і далі з ними чинять так само, як з ознаками форми: ступені, що збирають багато детекторів одного напрямку в різних місцях, відповідають на рух великого предмета незалежно від того, де він. Це та сама пара І та АБО~\eqref{eq:scstage}, тільки ознака не «край», а «край, що зсунувся за час $d$».

Відео ще й передбачуване: наступний кадр майже такий самий, як попередній. За гіпотезою \emph{передбачувального кодування} верхні ступені надсилають нижнім передбачення, а вгору йде головно помилка --- те, чого передбачення не врахувало~\cite{RaoBallard1999}. Це та сама економія імпульсів, що в твердженні~\ref{prop:economy}: платити за новини, а не за те, що й так відомо. Окремі спостереження з цією гіпотезою узгоджуються, але як загальну будову конвеєра її не доведено.

\section{Мозок і згорткова мережа}

Наскільки далеко сягає схожість? Для швидкого розпізнавання одного предмета вона справді глибока~\cite{Yamins2014}. Але в мозку є й те, чого немає в простій згортковій мережі.

\emph{Зворотні зв'язки.} Конвеєр мозку не лише прямий: у кожного ступеня є зв'язки назад і вбік. Для складних картинок --- з перекриттям, за поганого світла --- відповідь верхніх ступенів складається пізніше, і ці пізні відповіді краще описують мережі зі зворотними зв'язками, ніж прямі~\cite{Kar2019}. Один прохід розв'язує легкі випадки, складні потребують кількох кіл.

\emph{Стійкість.} Штучну мережу можна обдурити непомітною для ока підробкою пікселів~\cite{Szegedy2014}: зміна, якої людина не бачить, перетворює «панду» на «гібона»~\cite{Goodfellow2015}. Зір людини до таких підробок набагато стійкіший, але не невразливий: картинки, що обманюють одразу багато мереж, за короткого показу зсувають і вибір людини~\cite{Elsayed2018}. Чому стійкість так відрізняється --- питання відкрите; кандидати --- шум, ділення на загальну активність і зворотні зв'язки.

\emph{Учитель.} Мережі потрібні мільйони підписаних прикладів, мозку --- здебільшого відео без підписів~\eqref{eq:traceinvariance} і трохи підказок від \nt{DOPA} про те, що важливо.

\emph{Енергія.} Увесь мозок --- близько $20$ ватів (твердження~\ref{prop:economy}), а розпізнавання займає лише частину цього; навчена згорткова мережа на графічному прискорювачі витрачає сотні ватів.

\section{Ілюзії: де машина помиляється і чому}
\label{sec:illusions}

Інженер, який хоче зрозуміти чужу схему, подає на неї незвичні сигнали й дивиться, де вона помиляється. Для зору такі сигнали давно придумано --- це зорові ілюзії. Ілюзія --- не поломка. Кожна вказує на певний механізм машини, який у звичайних умовах працює правильно, а на спеціально дібраній картинці себе виказує.

\emph{Розумні очікування.} Вхід шумний, і машина доповнює його тим, що зазвичай буває у світі. Повільні рухи трапляються частіше за швидкі; коли вхід ненадійний --- наприклад, картинка малоконтрастна, --- оцінка швидкості зсувається до повільної, і блідий предмет здається рухомим повільніше за яскравий~\cite{Weiss2002}. Так само пояснюють багато ілюзій яскравості: ми бачимо не яскравість пікселя, а те, якій поверхні така яскравість найчастіше відповідала в минулому~\cite{Purves2011}. Фотографія сукні, яку одні бачать біло-золотою, а інші синьо-чорною, --- той самий механізм: картинка двозначна, і люди розходяться в тому, яким вони неусвідомлено вважають освітлення~\cite{LaferSousa2015,Wallisch2017}. Відмінності між людьми виміряно; те, що мозок тут поєднує вхід з очікуванням за правилами теорії ймовірностей, --- добре обґрунтована гіпотеза.

\emph{Регулювання підсилення.} Подивіться хвилину на водоспад, потім на нерухоме каміння: каміння попливе вгору. Канали «вниз» і «вгору» --- пара проводів, як у~\eqref{eq:dualrail}; регулювання підсилення~\eqref{eq:nakarushton} приглушило втомлений канал, і рівновага пари зсунулася. Ілюзії нахилу й контрасту, в яких оточення змінює вигляд центру, пояснюються діленням на активність сусідів~\cite{Schwartz2009}, як у розділі~\ref{sec:gaba}. Є й повчальний приклад обережності: темні плями на перехрестях білих смуг ґратки довго пояснювали простим гальмуванням сусідів, але досліди зі зміненою ґраткою показали, що це пояснення не працює~\cite{SchillerCarvey2005}.

\emph{Компенсація затримок.} Шлях від ока до верхніх ступенів займає десятки мілісекунд, і рухомий предмет за цей час устигає зсунутися. Якщо поруч із ним спалахне нерухома точка, предмет здається попереду спалаху, хоча вони збігалися~\cite{Nijhawan1994}. За одним із пояснень машина продовжує рух уперед, щоб компенсувати затримку, --- те саме, що в розділі~\ref{sec:muscles}: за запізнілого зворотного зв'язку доводиться передбачати. Пояснень у цієї ілюзії кілька, і суперечку не розв'язано.

\emph{Помилка в коді часу.} Нерухома картинка з кілець із різкими колірними переходами здається обертовою. Контрастні ділянки обробляються швидше за бліді, і різниця затримок~\eqref{eq:latency} читається детекторами напрямку як рух~\cite{Conway2005}. Запускають ілюзію дрібні мимовільні стрибки ока й кліпання: кожен стрибок оновлює вхід, і детектори знову бачать «рух»~\cite{OteroMillan2012}. Це код часом із розділу~\ref{sec:code}, прочитаний неправильно.

\emph{Дві картинки в одній.} Каркас куба, що дивиться на нас то однією гранню, то іншою, або різні картинки у двох очах: сприйняття перескакує кожні кілька секунд. Найкращі моделі --- взаємне гальмування двох тлумачень, повільна втома й шум~\cite{MorenoBote2007}, тобто та сама схема~\eqref{eq:halfcentre}, яка в розділі~\ref{sec:muscles} породжувала ритм кроку. Час зміни задають шум і втома; за моделлю~\cite{MorenoBote2007} головну роль відіграє шум, а втома лише допомагає.

А штучні мережі? Мережа, навчена передбачати наступний кадр відео, бачить обертання на тих самих нерухомих кільцях і не бачить його на контрольних картинках~\cite{Watanabe2018}. Мережі, навчені очищати зображення від шуму й підвищувати їхню різкість, піддаються тим самим ілюзіям кольору та яскравості, що й люди~\cite{GomezVilla2020}. Отже, частина ілюзій --- наслідок самої задачі, якої вчиться машина, а не особливостей нервової тканини. Які ілюзії властиві лише мозкові, --- добра перевірка для будь-якої моделі зору.

\begin{keyblock}
\begin{proposition}[Ілюзії як відбитки механізмів]
\label{prop:illusions}
Ілюзія виникає, коли механізм, корисний у звичайному світі, отримує незвичний вхід: очікування перемагає ненадійний вхід, регулювання підсилення зсуває пару каналів, компенсація затримки виносить предмет уперед, різниця затримок читається як рух, взаємне гальмування із шумом і втомою перескакує. Кожна ілюзія --- випробувальний сигнал, що вказує на свій механізм. Самі ілюзії виміряно; їхні пояснення --- від добре обґрунтованих до спірних.
\end{proposition}
\end{keyblock}

\section{Підсумок}

\begin{table}[t]
\centering
\footnotesize
\setlength{\tabcolsep}{4pt}
\renewcommand{\arraystretch}{1.15}
\begin{tabular}{>{\raggedright}p{3.0cm}>{\raggedright}p{4.6cm}>{\raggedright\arraybackslash}p{5.2cm}}
\toprule
Що робить машина & Як & Що відомо \\
\midrule
упізнає за $150\ms$ & один прохід десятком ступенів & швидкість виміряно \\
робить відповідь вибірковою & І за частинами, ступінь $S$~\eqref{eq:scstage} & виміряно на перших ступенях \\
робить відповідь інваріантною & АБО за положеннями, ступінь $C$; гальмування й ділення & виміряно на перших ступенях; для верхніх --- спрощення \\
видає відповідь & частоти близько сотні нейронів останнього ступеня, лінійне читання & виміряно \\
вчиться без підписів & правило Гебба зі слідом~\eqref{eq:traceinvariance} на неперервному відео & підміна предмета змінює відповіді --- виміряно; як головне правило --- гіпотеза \\
бачить рух & детектори напрямку, далі та сама пара І/АБО & виміряно \\
заощаджує на передбачуваному & угору йде помилка передбачення & гіпотеза \\
розв'язує складні випадки & зворотні зв'язки, кілька кіл & пізні відповіді й роль зворотних зв'язків виміряно \\
помиляється передбачувано & ілюзії: очікування, регулювання підсилення, затримки, взаємне гальмування із шумом і втомою & ілюзії виміряно; пояснення --- від обґрунтованих до спірних \\
\bottomrule
\end{tabular}
\caption{Як машина розпізнає зображення і відео.}
\label{tab:recognition}
\end{table}

Розпізнавання зібрано з деталей, які в нас уже були (таблиця~\ref{tab:recognition}): поріг дає І за частинами, взаємне гальмування --- АБО за положеннями, ділення --- незалежність від яскравості, правило Гебба зі слідом --- учителя, якого немає. Інженери взяли в зору чергування вибірковості й інваріантності та збудували на ньому згорткові мережі; навзаєм мозок поки не віддав головного --- вміння вчитися на відео без підписів майже без енергії.

\section{Задачі}

\begin{exercise}[\lvlA]
\label{ex:12:1}
\emph{Частота на одному ступені.} Ступінь конвеєра триває $15\ms$, нейрони працюють із частотою $20$ імпульсів за секунду, шуми корельовані з $c=0.1$~\eqref{eq:correlated}. Яка відносна похибка частоти групи зі ста нейронів і її межа для як завгодно великої групи? Скільки рівнів частоти розрізненні на ступені?
\end{exercise}

\begin{exercise}[\lvlA]
\label{ex:12:2}
\emph{Енергія одного впізнавання.} За один прохід за $150\ms$ десять ступенів по $10^7$ нейронів видають не більше одного імпульсу ціною $10^{-10}$~Дж. (а)~Яку частку енергії всього мозку за ці $150\ms$ коштує впізнавання? (б)~У скільки разів більше витрачає прискорювач потужністю $300$~Вт, що впізнає картинку за $10\ms$?
\end{exercise}

\begin{exercise}[\lvlB]
\label{ex:12:3}
\emph{Пороги ступеня $S$.} «Сітківка» з $24$ точок, фігури $A=11011$ і $B=10101$, зразок~\eqref{eq:scstage} дорівнює $+1$ на одиницях фігури й $-1$ на нулях. (а)~За яких порогів $S$-нейрон прощає одну перевернуту точку, але мовчить на чужій фігурі? (б)~Скільки $S$-нейронів потрібно для десяти ознак $5\times5$ на картинці $100\times100$?
\end{exercise}

\begin{exercise}[\lvlB]
\label{ex:12:4}
\emph{Скільки триває одна картинка з двох.} У схемі перескоків~\eqref{eq:halfcentre} при $\tau\ll\tau_a$, доки перемагає група~1, $r_2=0$ і $r_1=I-a_1$. (а)~Знайдіть тривалість перемоги при $I=1$, $\beta=2$, $b=2$, $\tau_a=0.4\,\text{с}$. (б)~Яке $\tau_a$ дає зміну картинки кожні три секунди?
\end{exercise}

\begin{exercise}[\lvlB]
\label{ex:12:5}
\emph{Моделювання: ціна інваріантності.} Функцією \texttt{two\_stage} з \texttt{models/\allowbreak recognition\_models.py} ($4000$ спроб) знайдіть, за якої ймовірності перевертання точки $p$ пара ступенів $S$, $C$ при $N=24$ помиляється в кожному четвертому випадку. Як змінюється частка правильних відповідей при $p=0.1$ від $N=8$ до $N=192$?
\end{exercise}

\begin{exercise}[\lvlB]
\label{ex:12:6}
\emph{Моделювання: вікно сліду.} Десять прогонів \texttt{trace\_learning} з \texttt{models/\allowbreak recognition\_models.py} з паузою $0.3\,\text{с}$: знайдіть, за яких $\tau_{\text{сл}}$ від $5\ms$ до $2\,\text{с}$ усі прогони вивчають інваріантність без помилок. Звідки в цього вікна верхня межа?
\end{exercise}

\begin{exercise}[\lvlB]
\label{ex:12:7}
\emph{Рух у будь-якому місці.} На сусідніх точках ряду з кроком $0.1^\circ$ стоять детектори напрямку розділу~\ref{sec:gaba}: гальмування від $A$ із затримкою $d$ гасить збудження від $B$, якщо вони розійшлися не більш ніж на $5\ms$; над ними --- ступінь $C$. Які затримки потрібні, щоб $C$ мовчав за зворотного руху зі швидкістю від $5$ до $20^\circ$ за секунду?
\end{exercise}

\begin{exercise}[\lvlC]
\label{ex:12:8}
\emph{Підробка пікселів.} Скільки перевернутих точок потрібно, щоб змінити відповідь моделі \texttt{two\_stage} ($N=24$) на чистій фігурі? А класифікатора «понад $3.5$ одиниці на вході --- це $A$», теж інваріантного до зсуву? Яка його частка правильних відповідей при $p=0.1$ проти $0.86$ у пари ступенів?
\end{exercise}
